% GENERATED FILE. Edit canonical lessons, then run npm run generate:books.
% canonical-source-hash: fnv1a64:5290dee7264b99c8
% canonical-lessons: TE-C321-situ, TE-C321-tikkettu-ivvandi, TE-C321-ausadham, TE-S177-letter-au, TE-C321-poyindi, TE-R321-at-the-station

\chapter{Tickets, Seats and Trouble}
\label{ch:te-tickets-seats-and-trouble}
\hlchaptermodality{321}

\begin{chapteropening}
\textbf{By the end of this chapter:} \emph{I can buy a ticket, ask for a seat and when a train leaves, and report a missed bus or a lost bag.}

Complete TE-R321-at-the-station: buy a ticket, ask for a seat and a time, report what went wrong, and read the new letter \te{ఔ}.
\end{chapteropening}

\section[\texorpdfstring{\te{సీటు} (sīṭu)}{సీటు (sīṭu)}]{\te{సీటు} (sīṭu) — a seat}

\label{lesson:TE-C321-situ}



\begin{quote}
Before the new one: say \textquotedblleft{}a bus\textquotedblright{} (\emph{bassu}) and \textquotedblleft{}a train\textquotedblright{} (\emph{railu}). Then ask \textquotedblleft{}how do I get to the station?\textquotedblright{} (\emph{sṭēṣanuki elā veḷḷāli?})
\end{quote}

\subsection*{You'll want to know: \te{సీటు}}
\textbf{\te{సీటు}} — \emph{sīṭu} — \textquotedblleft{}a seat\textquotedblright{}, on a bus or a train, or in a hall.

\begin{quote}
\textbf{\te{ఈ} \te{సీటు} \te{ఖాళీగా} \te{ఉందా}?}
\end{quote}

\begin{quote}
\emph{ī sīṭu khāḷīgā undā?} — \textquotedblleft{}Is this seat free?\textquotedblright{}
\end{quote}

\begin{quote}
\textbf{\te{నాకు} \te{కిటికీ} \te{దగ్గర} \te{సీటు} \te{కావాలి}.}
\end{quote}

\begin{quote}
\emph{nāku kiṭikī daggara sīṭu kāvāli.} — \textquotedblleft{}I would like a seat by the window.\textquotedblright{}
\end{quote}

\textbf{\te{ఖాళీగా}} (\emph{khāḷīgā}, \textquotedblleft{}empty, free\textquotedblright{}) with \textbf{\te{ఉందా}?}, \textquotedblleft{}is it?\textquotedblright{}: the \textbf{-\te{ఆ}} question again. Two seats are \textbf{\te{రెండు} \te{సీట్లు}} (\emph{reṇḍu sīṭlu}), the way two tickets are \textbf{\te{టిక్కెట్లు}}.

\begin{cousinweb}
\textbf{\te{సీటు}} (\emph{sīṭu}) is a loan from English \emph{seat}, with the \textbf{-\te{ఉ}} that Telugu gives most borrowed nouns, as in \textbf{\te{బస్సు}} and \textbf{\te{ఫోను}}.
\end{cousinweb}

\subsection*{Your turn}
Say these aloud:

\begin{itemize}
\raggedright
  \item \emph{sīṭu}
  \item \textquotedblleft{}Is this seat free?\textquotedblright{} — \emph{ī sīṭu khāḷīgā undā?}
  \item \textquotedblleft{}I would like a seat by the window\textquotedblright{} — \emph{nāku kiṭikī daggara sīṭu kāvāli}
\end{itemize}

\begin{tcolorbox}[breakable,colback=teal!4,colframe=teal!35!black,title={Before you move on}]
What does \textbf{\te{సీటు}} mean? (\textbf{A seat.}) \textquotedblleft{}Is this seat free?\textquotedblright{} (\textbf{\te{ఈ} \te{సీటు} \te{ఖాళీగా} \te{ఉందా}?}) \textquotedblleft{}Two seats\textquotedblright{}? (\textbf{\te{రెండు} \te{సీట్లు}}.)
\end{tcolorbox}
\section[\texorpdfstring{\te{ఒక} \te{టిక్కెట్టు} \te{ఇవ్వండి} (oka ṭikkeṭṭu ivvaṇḍi)}{ఒక టిక్కెట్టు ఇవ్వండి (oka ṭikkeṭṭu ivvaṇḍi)}]{\te{ఒక} \te{టిక్కెట్టు} \te{ఇవ్వండి} (oka ṭikkeṭṭu ivvaṇḍi) — at the ticket counter}

\label{lesson:TE-C321-tikkettu-ivvandi}



\begin{quote}
Say \textquotedblleft{}a ticket\textquotedblright{} (\emph{ṭikkeṭṭu}) and \textquotedblleft{}give\textquotedblright{}, politely (\emph{ivvaṇḍi}). Then \textquotedblleft{}is this seat free?\textquotedblright{} (\emph{ī sīṭu khāḷīgā undā?})
\end{quote}

\subsection*{You'll want to know: buying a ticket}
\begin{quote}
\textbf{\te{విజయవాడకి} \te{రెండు} \te{టిక్కెట్లు} \te{ఇవ్వండి}.}
\end{quote}

\begin{quote}
\emph{vijayavāḍaki reṇḍu ṭikkeṭlu ivvaṇḍi.} — \textquotedblleft{}Two tickets to Vijayawada, please.\textquotedblright{}
\end{quote}

The place takes \textbf{-\te{కి}}, \textquotedblleft{}to\textquotedblright{}; the number comes before the ticket; \textbf{\te{ఇవ్వండి}}, \textquotedblleft{}give\textquotedblright{}, closes the request.

\begin{quote}
\textbf{\te{టిక్కెట్టు} \te{ధర} \te{ఎంత}?}
\end{quote}

\begin{quote}
\emph{ṭikkeṭṭu dhara enta?} — \textquotedblleft{}How much is the ticket?\textquotedblright{}
\end{quote}

\textbf{\te{ఎంత}} (\emph{enta}, \textquotedblleft{}how much\textquotedblright{}). For the way back:

\begin{quote}
\textbf{\te{నేను} \te{రేపు} \te{తిరిగి} \te{వస్తాను}.}
\end{quote}

\begin{quote}
\emph{nēnu rēpu tirigi vastānu.} — \textquotedblleft{}I will come back tomorrow.\textquotedblright{}
\end{quote}

\textbf{\te{తిరిగి}}, \textquotedblleft{}back\textquotedblright{}, is the one in \textbf{\te{తిరిగి} \te{ఇవ్వు}}, \textquotedblleft{}to give back\textquotedblright{}.

\subsection*{You'll want to know: when does it leave?}
\begin{quote}
\textbf{\te{రైలు} \te{ఎన్ని} \te{గంటలకు} \te{బయలుదేరుతుంది}?}
\end{quote}

\begin{quote}
\emph{railu enni gaṇṭalaku bayaludērutundi?} — \textquotedblleft{}At what time does the train leave?\textquotedblright{}
\end{quote}

\begin{quote}
\textbf{\te{ఆరు} \te{గంటలకు} \te{బయలుదేరుతుంది}.}
\end{quote}

\begin{quote}
\emph{āru gaṇṭalaku bayaludērutundi.} — \textquotedblleft{}It leaves at six.\textquotedblright{}
\end{quote}

\textbf{\te{ఎన్ని} \te{గంటలకు}} (\emph{enni gaṇṭalaku}, \textquotedblleft{}at how many o'clock\textquotedblright{}): times of day take \textbf{-\te{కు}}, \textquotedblleft{}at\textquotedblright{}: \textbf{\te{ఆరు} \te{గంటలకు}}, \textquotedblleft{}at six\textquotedblright{}.

\subsection*{Your turn}
Say these aloud:

\begin{itemize}
\raggedright
  \item \textquotedblleft{}Two tickets to Vijayawada, please\textquotedblright{} — \emph{vijayavāḍaki reṇḍu ṭikkeṭlu ivvaṇḍi}
  \item \textquotedblleft{}How much is the ticket?\textquotedblright{} — \emph{ṭikkeṭṭu dhara enta?}
  \item \textquotedblleft{}At what time does the train leave?\textquotedblright{} — \emph{railu enni gaṇṭalaku bayaludērutundi?}
  \item buy one ticket to a town you know, ask the time, and say you will come back tomorrow
\end{itemize}

\begin{tcolorbox}[breakable,colback=teal!4,colframe=teal!35!black,title={Before you move on}]
Ask for two tickets, then ask the price and the time. (\textbf{\te{రెండు} \te{టిక్కెట్లు} \te{ఇవ్వండి}}; \textbf{\te{టిక్కెట్టు} \te{ధర} \te{ఎంత}}; \textbf{\te{ఎన్ని} \te{గంటలకు} \te{బయలుదేరుతుంది}}.)
\end{tcolorbox}
\section[\texorpdfstring{\te{ఔషధం} (auṣadhaṁ)}{ఔషధం (auṣadhaṁ)}]{\te{ఔషధం} (auṣadhaṁ) — a remedy; medication}

\label{lesson:TE-C321-ausadham}



\begin{quote}
Before the new one: say \textquotedblleft{}medicine\textquotedblright{} (\emph{mandu}), \textquotedblleft{}a fever\textquotedblright{} (\emph{jvaraṁ}) and \textquotedblleft{}a doctor\textquotedblright{} (\emph{vaidyuḍu}).
\end{quote}

\subsection*{You'll want to know: \te{ఔషధం}}
\textbf{\te{ఔషధం}} — \emph{auṣadhaṁ} — \textquotedblleft{}a remedy; medication\textquotedblright{}.

\begin{quote}
\textbf{\te{ప్రయాణంలో} \te{నాకు} \te{జ్వరం} \te{వచ్చింది}. \te{అక్కడ} \te{ఒక} \te{వైద్యుడు} \te{నాకు} \te{ఔషధం} \te{ఇచ్చాడు}.}
\end{quote}

\begin{quote}
\emph{prayāṇaṁlō nāku jvaraṁ vaccindi. akkaḍa oka vaidyuḍu nāku auṣadhaṁ iccāḍu.} — \textquotedblleft{}On the trip I got a fever. A doctor there gave me a remedy.\textquotedblright{}
\end{quote}

\begin{quote}
\textbf{\te{నవ్వు} \te{మంచి} \te{ఔషధం}.}
\end{quote}

\begin{quote}
\emph{navvu mañci auṣadhaṁ.} — \textquotedblleft{}Laughter is a good remedy.\textquotedblright{}
\end{quote}

Everyday Telugu says \textbf{\te{మందు}}. \textbf{\te{ఔషధం}} is the careful, learned word: you will read it on the signs of chemists' shops (\textbf{\te{ఔషధాలు}}, \textquotedblleft{}medicines\textquotedblright{}), on labels, and in sentences about what heals in a wider sense, like the one about laughter.

\begin{cousinweb}
\textbf{\te{ఔషధం}} (\emph{auṣadhaṁ}) is borrowed from Sanskrit \emph{auṣadha}, which CDIAL describes as \textquotedblleft{}made of herbs\textquotedblright{} and as \textquotedblleft{}a herb\textquotedblright{} (CDIAL 2571), with Telugu's noun ending \textbf{-\te{ం}}. A remedy was first a plant.
\end{cousinweb}

\subsection*{Your turn}
Say these aloud:

\begin{itemize}
\raggedright
  \item \emph{auṣadhaṁ}, the \emph{au} gliding from \emph{a} toward \emph{u}
  \item \textquotedblleft{}A doctor gave me a remedy\textquotedblright{} — \emph{oka vaidyuḍu nāku auṣadhaṁ iccāḍu}
  \item \textquotedblleft{}Laughter is a good remedy\textquotedblright{} — \emph{navvu mañci auṣadhaṁ}
  \item the everyday word for medicine, then the learned one
\end{itemize}

\begin{tcolorbox}[breakable,colback=teal!4,colframe=teal!35!black,title={Before you move on}]
What does \textbf{\te{ఔషధం}} mean? (\textbf{A remedy; medication.}) Which word is the everyday one? (\textbf{\te{మందు}}.) Where will you read \textbf{\te{ఔషధం}}? (\textbf{On chemists' signs and labels.})
\end{tcolorbox}
\section[\texorpdfstring{\te{ఔ} (au)}{ఔ (au)}]{\te{ఔ} (au) — the letter that opens \te{ఔషధం}}

\label{lesson:TE-S177-letter-au}



\begin{quote}
Before the new one: write \textbf{\te{ఒ}} and \textbf{\te{ఓ}} once each. Then the sign \textbf{\te{ౌ}}, which you know from \textbf{\te{గౌరవం}}, \textquotedblleft{}respect\textquotedblright{}.

Now say \textbf{\te{ఔషధం}} (\emph{auṣadhaṁ}). This lesson takes its first letter out and puts it in your hand.
\end{quote}

\subsection*{Script you'll notice: \te{ఔ}}
\textbf{\te{ఔ}} — \emph{au}, at the start of a word: one vowel that glides from \emph{a} toward \emph{u}, the way the middle of English \emph{now} does.

After a consonant, \emph{au} is the sign \textbf{\te{ౌ}}: \textbf{\te{గౌరవం}}. At the start of a word it stands as a letter of its own, \textbf{\te{ఔ}}. Few words open with it, and most of them come from Sanskrit.

In print \textbf{\te{ఔ}} is built on the body of \textbf{\te{ఒ}}, with a long stroke laid along its top that runs out to the right and ends in a small loop.

\subsection*{Writing: \te{ఔ} — copy the printed shape}
Put your pen beside \textbf{\te{ఔ}} and copy it, slowly, and larger than it is printed. Write the body of \textbf{\te{ఒ}} the way you learned it, then copy the long top stroke and its loop.

\begin{quote}
This book does not yet say where the pen starts on this shape, or which way it travels. That is taught with real variation from school to school, and it waits here until it can be written down from a source. Copying the printed shape needs no such source, and it puts the shape in your hand in the meantime.
\end{quote}

\subsection*{Your turn}
\begin{itemize}
\raggedright
  \item \emph{Find:} \te{ఔ} at the start of \textbf{\te{ఔషధం}}
  \item \emph{Trace it:} \te{ఔ} three times, saying \emph{au} as you finish each one
  \item \emph{Write it:} \te{ఒ}, \te{ఓ}, then \te{ఔ}, from memory, naming each
  \item \emph{Write it:} \textbf{\te{ఔషధం}}
\end{itemize}

\begin{tcolorbox}[breakable,colback=teal!4,colframe=teal!35!black,title={Before you move on}]
Which letter is this — \te{ఔ}? (\textbf{au}, at the start of a word.) Which word did it come from? (\textbf{\te{ఔషధం}}, \emph{auṣadhaṁ}, \textquotedblleft{}a remedy\textquotedblright{}.) Which sign says \emph{au} after a consonant? (\textbf{\te{ౌ}}, as in \textbf{\te{గౌరవం}}.)
\end{tcolorbox}
\section[\texorpdfstring{\te{పోయింది} (pōyindi)}{పోయింది (pōyindi)}]{\te{పోయింది} (pōyindi) — when something goes wrong}

\label{lesson:TE-C321-poyindi}



\begin{quote}
Say \textquotedblleft{}At what time does the train leave?\textquotedblright{} (\emph{railu enni gaṇṭalaku bayaludērutundi?}) Then say \textquotedblleft{}a remedy\textquotedblright{}. (\emph{auṣadhaṁ}.)

\begin{itemize}
\raggedright
  \item \emph{Write it:} \textbf{\te{ఔ}}, once
\end{itemize}
\end{quote}

\subsection*{You'll want to know: the bus has gone}
\begin{quote}
\textbf{\te{అయ్యో}! \te{బస్సు} \te{వెళ్ళిపోయింది}.}
\end{quote}

\begin{quote}
\emph{ayyō! bassu veḷḷipōyindi.} — \textquotedblleft{}Oh no! The bus has gone.\textquotedblright{}
\end{quote}

\textbf{\te{అయ్యో}} (\emph{ayyō}, \textquotedblleft{}oh no\textquotedblright{}). \textbf{\te{వెళ్ళిపోయింది}} is \textbf{\te{వెళ్ళి}}, \textquotedblleft{}having gone\textquotedblright{}, with \textbf{\te{పోయింది}}, \textquotedblleft{}went away\textquotedblright{}: gone, and not coming back. Telugu says what the bus did; English says what you missed.

\begin{quote}
\textbf{\te{నేను} \te{బస్సు} \te{అందుకోలేకపోయాను}.}
\end{quote}

\begin{quote}
\emph{nēnu bassu andukōlēkapōyānu.} — \textquotedblleft{}I could not catch the bus.\textquotedblright{}
\end{quote}

\textbf{\te{అందుకోలేకపోయాను}} (\emph{andukōlēkapōyānu}, \textquotedblleft{}I could not catch\textquotedblright{}) is \textbf{\te{అందుకో}}, \textquotedblleft{}to catch\textquotedblright{}, with \textbf{-\te{లేకపోయాను}}, \textquotedblleft{}I could not\textquotedblright{}.

\subsection*{You'll want to know: it is lost}
\begin{quote}
\textbf{\te{నా} \te{సంచి} \te{పోయింది}.}
\end{quote}

\begin{quote}
\emph{nā sañci pōyindi.} — \textquotedblleft{}My bag is lost.\textquotedblright{}
\end{quote}

\textbf{\te{పోవు}}, \textquotedblleft{}to go\textquotedblright{}, says \textquotedblleft{}to be lost\textquotedblright{} too: the bag \textquotedblleft{}went\textquotedblright{}. Then say what was in it and when it went:

\begin{quote}
\textbf{\te{సంచిలో} \te{నా} \te{ఫోను}, \te{నా} \te{డబ్బు} \te{ఉన్నాయి}.}
\end{quote}

\begin{quote}
\emph{sañcilō nā phōnu, nā ḍabbu unnāyi.} — \textquotedblleft{}My phone and my money were in the bag.\textquotedblright{}
\end{quote}

\begin{quote}
\textbf{\te{నిన్న} \te{సాయంత్రం} \te{స్టేషనులో} \te{పోయింది}. \te{నేను} \te{పోలీసులకు} \te{ఫిర్యాదు} \te{ఇవ్వాలి}.}
\end{quote}

\begin{quote}
\emph{ninna sāyantraṁ sṭēṣanulō pōyindi. nēnu pōlīsulaku phiryādu ivvāli.} — \textquotedblleft{}It went yesterday evening, at the station. I must make a complaint to the police.\textquotedblright{}
\end{quote}

\textbf{\te{ఉన్నాయి}} (\emph{unnāyi}, \textquotedblleft{}they are, they were\textquotedblright{}) is the plural of \textbf{\te{ఉంది}}. A report answers three questions: what, where, when.

\subsection*{Your turn}
Say these aloud:

\begin{itemize}
\raggedright
  \item \textquotedblleft{}Oh no! The bus has gone.\textquotedblright{} — \emph{ayyō! bassu veḷḷipōyindi.}
  \item \textquotedblleft{}I could not catch the bus\textquotedblright{} — \emph{nēnu bassu andukōlēkapōyānu}
  \item \textquotedblleft{}My bag is lost\textquotedblright{} — \emph{nā sañci pōyindi}
  \item \textquotedblleft{}My phone and my money were in the bag\textquotedblright{} — \emph{sañcilō nā phōnu, nā ḍabbu unnāyi}
  \item report a lost phone: what, where and when, and that you must make a complaint
\end{itemize}

\begin{tcolorbox}[breakable,colback=teal!4,colframe=teal!35!black,title={Before you move on}]
\textquotedblleft{}The bus has gone\textquotedblright{}? (\textbf{\te{బస్సు} \te{వెళ్ళిపోయింది}}.) \textquotedblleft{}My bag is lost\textquotedblright{}? (\textbf{\te{నా} \te{సంచి} \te{పోయింది}}.) What three questions does a report answer? (\textbf{What, where, when.})
\end{tcolorbox}
\section[(review)]{Review — at the station}

\label{lesson:TE-R321-at-the-station}



\begin{quote}
Say \textquotedblleft{}a remedy\textquotedblright{}. (\emph{auṣadhaṁ}.) Then the three words whose first letters came before it. (\emph{ūru}, \emph{ōṭu}, \emph{īlōgā}.)

\begin{itemize}
\raggedright
  \item \emph{Write it:} \textbf{\te{ఔ}}, \textbf{\te{ఊ}}, \textbf{\te{ఓ}} and \textbf{\te{ఈ}}, once each from memory
\end{itemize}
\end{quote}

\subsection*{Your turn}
\begin{itemize}
\raggedright
  \item \emph{Say it:} One ticket to my home town, please. At what time does the bus leave?
  \item \emph{Say it:} Is this seat free? I would like a seat by the window.
  \item \emph{Say it:} Oh no, the bus has gone. How do I get to the station?
  \item \emph{Say it:} My bag is lost. My phone was in the bag.
  \item \emph{Write it:} \textbf{\te{ఔషధం}}, \textbf{\te{ఊరు}}, \textbf{\te{ఓటు}}, \textbf{\te{ఈలోగా}}
\end{itemize}

\begin{tcolorbox}[breakable,colback=teal!4,colframe=teal!35!black,title={Before you move on}]
\textquotedblleft{}A seat\textquotedblright{}? (\textbf{\te{సీటు}}.) \textquotedblleft{}At what time does it leave?\textquotedblright{} (\textbf{\te{ఎన్ని} \te{గంటలకు} \te{బయలుదేరుతుంది}?}) \textquotedblleft{}It has gone\textquotedblright{}? (\textbf{\te{వెళ్ళిపోయింది}}.) \textquotedblleft{}It is lost\textquotedblright{}? (\textbf{\te{పోయింది}}.) \textquotedblleft{}A remedy\textquotedblright{}? (\textbf{\te{ఔషధం}}, which opens with \textbf{\te{ఔ}}.)
\end{tcolorbox}
